% !PRINTED 125
\weberpage{125}

Die genaue Bestimmung dieser Konstanten $\varepsilon$, namentlich auch mit Rücksicht auf das Vorzeichen ist ein bekanntes wichtiges Problem, das eigentümliche Schwierigkeiten bietet, dessen Lösung aber für uns unerläßlich ist. Lösungen haben auf verschiedenen Wegen \gesperrt{Hermite}\footnote{Liouvilles Journal, Ser.\,II, T.\,III, 1858. Oeuvres de \gesperrt{Charles Hermite}, p.\,487.} und \gesperrt{Dedekind}\footnote{Erläuterungen zu Nr.\,XXVIII von \gesperrt{Riemanns Werken}, zweite Auflage und \glqq{}Über die elliptischen Modulfunktionen\grqq{}, Crelles Journal, Bd.\,83, S.\,265. Vgl. auch des Verfassers Abhandlung \glqq{}Zur Theorie der elliptischen Funktionen\grqq{}, Acta Mathematica, Bd.\,6, S.\,341\,ff.}, neuerdings auch \gesperrt{Mertens} und \gesperrt{Scheibner}\footnote{\gesperrt{Mertens}, Zur linearen Transformation der $\vartheta$-Reihen, Transactions of the American mathematical society. July 1901. — \gesperrt{Scheibner}, Zur linearen Transformation der Theta-Funktionen und elliptischen Modulfunktionen, Berichte der Sächs. Gesellschaft der Wissenschaften. Oktober 1906.} gegeben. Wir wollen hier einen Weg gehen, der den Vorzug großer Einfachheit hat, dafür freilich nicht eine Ableitung, sondern nur einen Beweis der fertigen Formel enthält.

Für zwei spezielle Fälle haben wir schon früher [§ 34, (4), (5)] diese Bestimmung ausgeführt und auf dies Ergebnis werden wir uns hier stützen. Es sind die Formeln:
\weq{5}{\eta(\omega \pm 1) = e^{\pm\frac{\pi i}{12}}\,\eta(\omega)}
und
\weq{6}{\eta\left(-\frac{1}{\omega}\right) = \sqrt{-i\omega}\ \eta(\omega),}
worin $\sqrt{-i\omega}$ \gesperrt{mit positivem reellem Teil} zu nehmen ist.

Wir setzen nun
\weq{7}{\frac{\eta\left(\dfrac{\gamma+\delta\omega}{\alpha+\beta\omega}\right)}{\eta(\omega)} = E\left(\begin{smallmatrix}\alpha, & \beta\\ \gamma, & \delta\end{smallmatrix}, \omega\right),}
und haben den Wert dieses Symbols zu bestimmen. Nach (5), (6), (7) haben wir:
\weq{8}{E\left(\begin{smallmatrix}-\alpha, & -\beta\\ -\gamma, & -\delta\end{smallmatrix}, \omega\right) = E\left(\begin{smallmatrix}\alpha, & \beta\\ \gamma, & \delta\end{smallmatrix}, \omega\right),}
\weq{9}{E\left(\begin{smallmatrix}1, & 0\\ 0, & 1\end{smallmatrix}, \omega\right) = 1,}
\weq{10}{E\left(\begin{smallmatrix}1, & 0\\ \pm 1, & 1\end{smallmatrix}, \omega\right) = e^{\pm\frac{\pi i}{12}},}
\weq{11}{E\left(\begin{smallmatrix}0, & 1\\ -1, & 0\end{smallmatrix}, \omega\right) = \sqrt{-i\omega}.}
